%% kapa: exam=1 round=27 year=2016 period=2 form=A subject=accounting subject_ko=회계학 seq=20 number=60
%% source: Q-Net 감정평가사 기출문제 (cst00309) artlSeq=5169402 file=extracted/1008/5169402/제27회 감정평가사 1차 문제/제27회 감정평가사 2교시문제(A형).pdf
%% conversion: kordoc 4.24.51 -> tools/gen1.py (auto); review=proofread-sonnet; flags=none
\begin{question}{60}
20×6년 1월 1일 (주)감평은 건물과 토지를 ￦2,000,000에 일괄구입하였다. 구입당시 건물과 토지의 공정가치는 각각 ￦960,000과 ￦1,440,000이었다. 건물의 내용연수는 7년, 잔존가치는 ￦100,000으로 추정하였으며 정액법으로 감가상각한다. 20×6년 12월 31일 건물과 토지에 관한 순공정가치와 사용가치는 다음과 같으며 회수가능액과 장부금액의 차이는 중요하고 손상징후가 있다고 판단된다.
\begin{framedtext}
{\centering\begin{tabular}{l@{\hspace{2em}}r@{\hspace{2em}}r}
 & \uline{\hspace{0.6em}순공정가치\hspace{0.6em}} & \uline{\hspace{0.6em}사용가치\hspace{0.6em}} \\
건물 & ￦600,000 & ￦670,000 \\
토지 & 1,150,000 & 1,000,000 \\
\end{tabular}\par}
\end{framedtext}
(주)감평이 20×6년도에 인식해야 할 손상차손은?
\begin{choices}
  \item ￦0
  \item ￦80,000
  \item ￦130,000
  \item ￦230,000
  \item ￦300,000
\end{choices}
\end{question}
\answer{2}
